% Nouveau EGA V (cote 25): archive cote 25, site page 77 (PDF page 78).
% typescript. Draft compared: none.
% AI-transcribed by Claude Opus 5.5 (claude-opus-5-5, adaptive thinking, effort high), Message Batches API, batch ega25b (msgbatch_01A7n5YuXubUPNK1zsepi8RG), 2026-11-13; not yet checked by eye.
\scanpage{25/77}
\mspage{\hand{72}}

une sous-variété linéaire donnée C  de codimension n.  A cet effet,
on aura besoin du

\uline{Lemme} 13.3.  Avec les notations précédentes pour \xout{xx} X,Y,Z  ($X\supset Y\supset Z$
sous schémas fermés de $P^r_k$ ,\inserted[typed]{de dimensions n,n-1 et \hand{$\leqslant$} n-2} X-Z\hand{,} \deleted{lisse de dim n en tous ses points},
Y-Z lisse\hand{\unsure{s}} \deleted{de dim n-1 en tous ses points;}\hand{,} Z de dimension \hand{$\leqslant$}n-2), et quitte
au besoin (lorsque k est de car p$>$0)  de remplacer l'immersion projecti-%
ve  i:X$\rightarrow$P$^r$  par  un ``multiple'' quelconque $\phi_n i$ ($n \hand{$\geqslant$} 2$)  comme dans le
n$^o$ 9,  il existe une sous-variété linéaire  $L_o$ de P$^r$ , de codimension
n-1, et ayant les vertus suivantes:

a)  $L_o \cap Z = \emptyset$ .

b)  $L_o \cap X$  est lisse de dimension 1.

c)  $L_o \cap Y$   est lisse de dimension  0.

(NB  k désigne un corps infini, pas la peine ici qu'il soit parfait).

Admettons  ce lemme  et montrons comment on en déduit  l'existence
d'un voisinage ouvert U de x, contenu dans X-Y, et admettant un morphis-%
me élémentaire U$\longrightarrow$V.  \deleted{En vertu du n$^o$ 7 appliqué à l'espace projectif}
\deleted{$L_o$ et au sous-schéma lisse $L_o\cap Y$ de celui-ci,} \over{i}{I}l existe une sous-%
variété linéaire C de $L_o$ , de codimension n dans P$^r$ i.e. de codimension
1 dans $L_o$ , \inserted{ne rencontrant pas l'ensemble fini \deleted{\illegible{1}} $(L_o\cap Y)\cup\{x\}$.} \deleted{telle que C $\cap$ X ( = C $\cap$ ($L_o\cap$ X) ) soit lisse sur k et de}
\deleted{dimension 0; de plus, on peut exiger que} l'on ait  $C \cap Y = \emptyset$ \inserted{et $x\notin C$,} \deleted{(même}
\deleted{argument).} Soit  T = X$\cap$C, de sorte que T est un sous-schéma de X,
étale sur k, non vide,\inserted{ne contenant pas x,} et \uline{disjoint} de Y.  Considérons d'autre part le
sous-schéma Q de Grass$^{n-1}$(P) correspondant aux sous-variétés linéaires
de P$^r$ contenant \unsure{C|E}, de sorte que  Q est un \xout{espace schéma} espace projectif
de dimension n-1,   en particulier il est lisse sur k de dim. n-1.
Alors $L_o$ correspond à un point \hand{$\xi$}$_o$ de Q(k).  Considérons d'autre part
(avec les notations générales introduites par ailleurs)  l'image inverse
$\underline{X}^{n-1}_Q$ de $\underline{X}^{n-1}$  par  l'immersion Q$\longrightarrow$ Grass$^{n-1}$ , et de même les images

\margin{77}
